% !TITLE 辛夷坞
% !AUTHOR 王维
% !DYNASTY 唐
% !GENRE 五言绝句
% !ORDER 43

\begin{poem}
木末芙蓉花\textsuperscript{1}，山中发红萼\textsuperscript{2}。
涧户\textsuperscript{3}寂无人，纷纷开且落。
\end{poem}

\begin{annotations}
\notetext{1}{木末芙蓉花：树梢上的芙蓉花，指辛夷（木兰）花。木末，树梢。辛夷花色红形似芙蓉。}
\notetext{2}{红萼：红色的花萼（花托），指花瓣。在一片绿山中，那一点红格外醒目。}
\notetext{3}{涧户：山涧的入口处，指偏僻的山谷。花儿在空寂无人的山谷中自开自落——不为任何人。}
\end{annotations}

\begin{appreciation}
辛夷坞是辋川别业附近的一处山谷，辛夷就是木兰花。这首诗写的是山中木兰花自开自落的景象。

木末芙蓉花，山中发红萼——树梢上的芙蓉花，在山中绽放出红色的花萼。木末就是树梢。萼（è）是花托，这里指花瓣。王维不说木兰，而说芙蓉花，因为辛夷花的颜色和形状都像芙蓉。红萼二字，写出了花的鲜艳和孤独：在一片绿色的山林中，那一点红色格外醒目，却无人欣赏。

涧户寂无人，纷纷开且落——山涧边寂静无人，花儿纷纷开放，又纷纷飘落。涧户指山涧的入口，这里指山谷。纷纷二字，写出了花开的繁盛和落花的无奈。开且落，三个字道尽了生命的全过程：从绽放到凋零，不过是一瞬间的事。

这首诗表面上是写景，实际上是在说一种生命的态度。花不为任何人而开，也不为任何人而落。它开，是因为春天到了；它落，是因为时候到了。没有观众，没有掌声，生命的完成本身就是意义。

六百年后，明代的王阳明在《传习录》中说了一段著名的话：你未看此花时，此花与汝心同归于寂；你来看此花时，则此花颜色一时明白起来。这话常被用来阐释心学的主观唯心主义，但若回到王维的诗中，你会发现另一种意韵：花自在开落，不因人看或不看而改变。它的明白，是它自己的明白；它的寂灭，是它自己的寂灭。人看或不看，不过是多了一个旁观者罢了。

王维晚年隐居辋川，写了很多这样的小诗。它们不宏大，不激烈，只是静静地记录着一个又一个瞬间。但正是这些瞬间，让我们看到了中国诗歌中最深邃的宁静。
\end{appreciation}
